
\documentclass[12pt,a4paper]{article}

\usepackage{fontspec}
\usepackage[english,slovak]{babel}

\usepackage{graphicx}
\usepackage{booktabs}
\usepackage{float}
\usepackage{placeins}
\usepackage{emoji}
\usepackage{hyperref}
\usepackage{enumitem}
\usepackage{amssymb}

\usepackage[
  a4paper,
  margin=2.5cm
]{geometry}

\usepackage{amsmath,amssymb}

\title{Report skupiny
  \emoji{folded-hands}
  \emoji{folded-hands}
  \emoji{folded-hands}
}
\author{Vladislav Voenkov, Petra Nagyová}
\date{}

\begin{document}

\maketitle


\section*{Úloha A}

Najprv potrebujeme diskretizovať účelovú funkciu spotreby paliva
\[
J(f) = \gamma \int_{0}^{T^p} \left\|f(t)\right\|_{2} dt =
\gamma \sum_{k=1}^K \left[ \int_{(k-1)h}^{kh} \left\|f(t)\right\|_{2} dt \right].
\]
Keďže predpokladáme, že počas jednej periódy \(k\) je ťahová sila konštantná
a rovná sa \(f_k\), funkciu vieme prepísať ako
\[
J(f) = \gamma \sum_{k=1}^K \left[ \int_{(k-1)h}^{kh} \left\|f_k\right\|_{2} dt \right] =
\gamma \sum_{k=1}^K h \left\|f_k\right\|_{2} =
\gamma h \sum_{k=1}^K \left\|f_k\right\|_{2}.
\]
Ak chceme nájsť len optimálne hodnoty \(f_k\), vieme vynechať kladné konštanty
\(\gamma\) a \(h\). Potom úlohu vieme zapísať nasledovne:
\[
\begin{alignedat}{3}
  \min_{p,v,f} \quad
  & \sum_{k=1}^{K} \|f_k\|_2 \\
  & p_1 - p^{init} &\quad&= 0 \\
  & v_1 - v^{init} &\quad&= 0 \\
  & p_{K+1}       &\quad&= 0 \\
  & v_{K+1}       &\quad&= 0 \\
  & v_k - v_{k+1} + \frac{h}{m} f_k - h g e_3 &\quad&= 0
  &\quad&\forall k \in \{1,\dots,K\} \\
  & p_k - p_{k+1} + \frac{h}{2} v_k + \frac{h}{2} v_{k+1} &\quad&= 0
  &\quad&\forall k \in \{1,\dots,K\}\\
  & \|f_k\|_2 - F^{\max} &\quad&\le 0 &\quad&\forall k \in \{1,\dots,K\} \\
  & \alpha \left\|
    \begin{pmatrix}
      (p_k)_1 \\
      (p_k)_2
    \end{pmatrix}
    \right\|_2 - (p_k)_3
    &\quad&\le 0
    &\quad&\forall k \in \{1,\dots,K+1\}.
\end{alignedat} \quad (U)
\]

\begin{samepage}
  Ukážme, že táto úloha patrí medzi úlohy konvexného programovania.
  Účelová funkcia je konvexná, pretože je súčtom euklidovských noriem, ktoré sú
  konvexnými funkciami. Obmedzenia s rovnosťami sú afinné v rozhodovacích
  premenných. Funkcia
  \[
  \|f_k\|_2 - F^{\max}
  \]
  je konvexná, pretože vznikne posunutím konvexnej funkcie o konštantu. Funkcia
  \[
  \alpha \left\|
  \begin{pmatrix}
    (p_k)_1 \\
    (p_k)_2
  \end{pmatrix}
  \right\|_2 - (p_k)_3
  \]
  je konvexná, pretože je súčtom konvexnej a lineárnej funkcie.
\end{samepage}

Takže účelová funkcia je konvexná, obmedzenia s rovnosťami sú afinné,
obmedzenia s nerovnosťami sú konvexné, preto daná úloha patrí medzi úlohy
konvexného programovania. Konkrétne ide o úlohu kužeľového
programovania druhého rádu (\emph{second-order cone program, SOCP})
\footnote{\url{https://en.wikipedia.org/wiki/Second-order_cone_programming}}.

Zavedením pomocných premenných \(t_k\) môžeme prepísať účelovú funkciu do
lineárneho tvaru
\[
\min_{p,v,f,t} \quad \sum_{k=1}^{K} t_k
\]
s dodatočnými obmedzeniami
\[
\|f_k\|_2 - t_k \le 0 \quad \forall k \in \{1, \dots, K\}.
\]
Potom výsledná úloha má lineárnu účelovú funkciu, obmedzenia s rovnosťami
sú afinné a obmedzenia s nerovnosťami majú tvar nerovností druhého rádu, t. j.
\[
\left\| A_k x + b_k \right\|_2 \le c_k^\intercal x + d_k,
\]
čiže ekvivalentná úloha patrí do triedy SOCP.



\section*{Úloha B}

Na rozdiel od úlohy A, kde sme pri pevnom počte periód $K$
minimalizovali spotrebu paliva, teraz hľadáme najmenší
počet periód $K_{\min}$, pre ktorý existuje prípustný
pristávací manéver.
Na tento účel využijeme obmedzenia úlohy $(U)$ a ukážeme,
že prípustnosť je vzhľadom na počet periód monotónna.

\subsection*{Nutná podmienka prípustnosti}

Najprv ukážeme, že existencia prípustného pristávacieho
manévra vyžaduje $mg \le F^{\max}$.

\textbf{Tvrdenie:}
Ak $mg > F^{\max}$, neexistuje žiadny prípustný
pristávací manéver pre ľubovoľné $K \ge 1$.

\textit{Dôkaz.}
Predpokladajme sporom, že $mg > F^{\max}$ a zároveň
existuje prípustný manéver pre nejaké $K \ge 1$.
Z obmedzenia ťahovej sily a vlastností noriem dostávame
\[
(f_k)_3 \le \|f_k\|_2 \le F^{\max} < mg,
\qquad \forall k \in \{1,\dots,K\}.
\]
Z dynamickej rovnice pre rýchlosť potom vyplýva
\[
(v_{k+1})_3
= (v_k)_3 + \frac{h}{m}(f_k)_3 - hg
< (v_k)_3.
\]
Zvislá rýchlosť teda striktne klesá.
Keďže $(v_k)_3 > 0$ pre akekolvek $k \in \{1,\dots,K\}$
a $(v_{K+1})_3 = 0$, z dynamickej rovnice pre polohu
dostávame
\[
(p_{k+1})_3 - (p_k)_3
= \frac{h}{2}\big((v_k)_3 + (v_{k+1})_3\big) > 0
\]
pre všetky $k \in \{1,\dots,K\}$.

To znamená, že zvislá poloha lode sa počas každej
periódy striktne zväčšuje, teda
\[
(p_1)_3 < (p_2)_3 < \dots < (p_{K+1})_3.
\]

Z koncovej podmienky však vieme, že loď musí
pristáť v počiatku súradnicového systému,
čiže $(p_{K+1})_3 = 0$. Z predchádzajúcej
nerovnosti preto vyplýva
\[
(p_1)_3 < (p_{K+1})_3 = 0,
\]
teda počiatočná výška lode musí byť záporná.

To však nie je možné, pretože podľa kužeľového
obmedzenia musí platiť
\[
(p_1)_3 \ge
\alpha \left\|
\begin{pmatrix}
(p_1)_1\\
(p_1)_2
\end{pmatrix}
\right\|_2 \ge 0.
\]

Dostávame teda spor:
\[
\boxed{(p_1)_3 < 0 \quad \land \quad (p_1)_3 \ge 0.}
\]

Preto pri $mg > F^{\max}$ neexistuje žiadny
prípustný pristávací manéver.
\hfill $\blacksquare$


\subsection*{Monotónnosť prípustnosti}

Predpokladajme, že existuje prípustný pristávací
manéver $(p,v,f)$ pre $K$ periód. Ukážeme, že ho
vieme rozšíriť na prípustný manéver pre $K+1$ periód.

\textbf{Idea:}
Po pristátí loď zostane v pokoji ešte jednu
periódu dĺžky $h$.

Pôvodný manéver ponecháme nezmenený a zadefinujeme
nové premenné
\[
p_{K+2}=0,\qquad
v_{K+2}=0,\qquad
f_{K+1}=mg e_3.
\]

Keďže $p_{K+1}=v_{K+1}=0$, dynamické rovnice
pre novú periódu sú splnené:
\[
\begin{aligned}
v_{K+2}
&=v_{K+1}+\frac{h}{m}f_{K+1}-hg e_3=0,\\
p_{K+2}
&=p_{K+1}+\frac{h}{2}(v_{K+1}+v_{K+2})=0.
\end{aligned}
\]

Kužeľové obmedzenie je v novej polohe $p_{K+2}=0$
triviálne splnené. Zároveň platí
\[
\|f_{K+1}\|_2 = mg \le F^{\max},
\]
pričom posledná nerovnosť vyplýva
z predchádzajúceho tvrdenia, keďže pôvodný
manéver je prípustný.

Rozšírený manéver teda spĺňa všetky obmedzenia
úlohy $(U)$ pre $K+1$ periód.

\textbf{Záver:}
Ak existuje prípustný pristávací manéver pre
$K$ periód, potom existuje aj pre $K+1$ periód,
a teda aj pre ľubovoľný väčší počet periód.

\subsection*{Hľadanie minimálneho času pristátia}

Z predchádzajúceho výsledku vyplýva, že
prípustnosť je vzhľadom na počet periód $K$
monotónne neklesajúca.

Ak existuje minimálny prípustný počet periód
$K_{\min}$, potom platí
\[
\begin{cases}
K<K_{\min}
&\Rightarrow \text{úloha je neprípustná},\\
K\ge K_{\min}
&\Rightarrow \text{úloha je prípustná}.
\end{cases}
\]

Pre pevné $K$ môžeme prípustnosť overiť riešením úlohy
\[
\min_{p,v,f}\quad 0
\]
so všetkými obmedzeniami úlohy $(U)$.

Keďže v úlohe A sme ukázali, že tieto obmedzenia
definujú konvexnú množinu, a účelová funkcia
je konštantná, ide o konvexnú optimalizačnú úlohu,
konkrétne o SOCP.

Vďaka monotónnosti môžeme $K_{\min}$ nájsť
kombináciou exponenciálneho a binárneho vyhľadávania.

\textbf{1. fáza: Exponenciálne vyhľadávanie}

Najprv potrebujeme nájsť prípustnú hornú hranicu
$K_R$.

\begin{enumerate}[label=\arabic*.]
  \item Zvolíme maximálny skúmaný počet periód
  $K_{\max}\ge1$ a položíme
  \[
  K_L:=0,\qquad K_R:=1.
  \]

  \item Overíme prípustnosť úlohy pre $K_R$ periód.

  \item Ak je úloha prípustná, prejdeme
  na druhú fázu.

  \item Ak $K_R=K_{\max}$, ukončíme algoritmus.
  V rozsahu $1\le K\le K_{\max}$ neexistuje
  prípustný pristávací manéver.

  \item Inak nastavíme
  \[
  K_L:=K_R,\qquad
  K_R:=\min(2K_R,K_{\max})
  \]
  a vrátime sa na krok 2.
\end{enumerate}

Po úspešnom ukončení prvej fázy máme prípustnú
hornú hranicu $K_R$ a dolnú hranicu $K_L$,
pričom platí $K_L<K_{\min}\le K_R$.
Hodnota $K_L=0$ predstavuje hranicu pod
skúmaným rozsahom $K\ge1$.

\textbf{2. fáza: Binárne vyhľadávanie}

V intervale $(K_L,K_R]$ hľadáme najmenší
prípustný počet periód.

\begin{enumerate}[label=\arabic*.]
  \item Pokiaľ $K_R-K_L>1$, vykonáme
  nasledujúce kroky:

  \begin{enumerate}[label=\arabic*.]
    \item Položíme
    \[
    K:=\left\lfloor\frac{K_L+K_R}{2}\right\rfloor.
    \]

    \item Pomocou solvera overíme prípustnosť
    úlohy pre $K$ periód.

    \item Ak je úloha prípustná, nastavíme
    $K_R:=K$.

    \item Inak nastavíme $K_L:=K$.
  \end{enumerate}

  \item Keď $K_R-K_L=1$, dostávame
  \[
  K_{\min}=K_R,\qquad
  T_{\min}^{p}=K_{\min}h.
  \]
\end{enumerate}

\end{document}
